% Part: conditional-logics
% Chapter: minimal-change-semantics
% Section: sphere-models

\documentclass[../../../include/open-logic-section]{subfiles}

\begin{document}

\olfileid[tr]{con}{min}{sph}

\olsection{Küre Modelleri}

Karşıolgusallara biçimsel semantik vermenin bir yolu, Lewis'in biçimsel
olmayan açıklamasını matematiksel bir yapıya dönüştürmektir. Bu durumda bir
$w$ dünyasının çevresindeki küreler dünya kümeleridir. Küreler iç içe
olduğundan, $w$ çevresindeki dünya kümelerinin alt küme bağıntısıyla doğrusal
olarak sıralanması gerekir.

\begin{defn}
  Bir \emph{küre modeli}, $W$'nin boş olmayan bir dünyalar kümesi,
  $V\colon\PVar\to\Pow{W}$'nin bir değerleme ve
  $O\colon W\to\Pow{\Pow{W}}$'nin her $w$ dünyasına bir
  \emph{küreler sistemi}~$O_w$ atadığı $\mModel{M}=\tuple{W,O,V}$
  üçlüsüdür. Her $w$ için $O_w$ bir dünya kümeleri kümesidir ve şunları
  sağlamalıdır:
  \begin{enumerate}
  \item $O_w$, $w$ üzerinde \emph{merkezlidir}: $\{w\}\in O_w$.
  \item $O_w$ \emph{iç içedir}: $S_1,S_2\in O_w$ olduğunda
    $S_1\subseteq S_2$ ya da $S_2\subseteq S_1$; yani $O_w$,
    $\subseteq$ ile doğrusal sıralıdır.
  \item $O_w$, boş olmayan birleşimler altında kapalıdır.
  \item $O_w$, boş olmayan kesişimler altında kapalıdır.
  \end{enumerate}
\end{defn}

$O_w$'nin ardındaki sezgi, $w$'nin ``çevresindeki'' dünyaların $w$'den ne
kadar uzakta olduklarına göre katmanlara ayrılmasıdır. En içteki küre yalnızca
$w$'nin kendisi, yani $\{w\}$ kümesidir: $w$, başka her küredeki dünyalardan
$w$'ye daha yakındır. $S\subsetneq S'$ ise $S'\setminus S$ içindeki dünyalar
$w$'den, $S$ içindeki dünyalara göre daha uzaktır: $S'\setminus S$, $S$ ile
$S'$'nin dışındaki dünyalar arasındaki ``katmandır.'' Özellikle kürelerin dış
yüzeylerinin içindeki bütün dünyaları içerdiğini düşünmeliyiz; küreler yalnızca
tek tek katmanlardan ibaret değildir.

\begin{figure}
\begin{center}
\begin{tikzpicture}[layerwidth=1.5,scale=.6]\tiny
  \spheresystem{3}
  \propositionintersect{0}{3}{60}{6}
  \path (0,0) node[world] {$w$};
  \spherepos{-30}{2}{node[world] {$w_2$}}
  \spherepos{220}{2}{node[world] {$w_3$}}
  \spherepos{90}{2}{node[world] {$w_1$}}
  \spherepos[shift={(.1,0)}]{10}{3}{node[world] {$w_5$}}
  \spherepos[shift={(.1,0)}]{-10}{3}{node[world] {$w_6$}}
  \spherepos{225}{3}{node[world] {$w_4$}}
  \spherepos{20}{4}{node[world] {$w_7$}}
  \path (5,0) node[left] {$p$};
\end{tikzpicture}
\caption{Bir küre modelinin şeması}
\ollabel{fig:sphere-model}
\end{center}
\end{figure}

\olref{fig:sphere-model}'deki şema, $W=\{w,w_1,\dots,w_7\}$ ve
$V(p)=\{w_5,w_6,w_7\}$ olan küre modeline karşılık gelir. En içteki küre
$S_1=\{w\}$'dir. $w$'ye en yakın dünyalar $w_1,w_2,w_3$ olduğundan bir
sonraki büyük küre $S_2=\{w,w_1,w_2,w_3\}$'tür. Daha dışarıdaki dünyalar
$w_4,w_5,w_6$ olduğundan en dış küre $S_3=\{w,w_1,\dots,w_6\}$'dır.
$w$ çevresindeki küreler sistemi $O_w=\{S_1,S_2,S_3\}$'tür. $w_7$ dünyası
$w$ çevresindeki hiçbir kürede değildir. $p$'nin doğru olduğu en yakın
dünyalar $w_5$ ile~$w_6$'dır; dolayısıyla $p$'yi kabul eden en küçük küre
$S_3$'tür.

Bir küre modeli~$\mModel M$ içinde $w$ dünyasında bir $!A$ !!{formula}{}ünün
sağlanmasını, $\mSat{M}{!A}[w]$'yi tanımlamak için modal !!{formula}s{}in
tanımını $!B\cif !C$ için bir koşul ekleyerek genişletiriz:

\begin{defn}
  $\mSat{M}{!B \cif !C}[w]$ ancak ve ancak aşağıdakilerden biri geçerlidir:
  \begin{enumerate}
  \item\ollabel{sphere-vac} Her $u\in\bigcup O_w$ için
    $\mSat/{M}{!B}[u]$; ya da
  \item\ollabel{sphere-nonvac} Bazı $S\in O_w$ için:
    \begin{enumerate}
      \item bazı $u\in S$ için $\mSat{M}{!B}[u]$ ve
      \item her $v\in S$ için ya $\mSat/{M}{!B}[v]$ ya da
        $\mSat{M}{!C}[v]$.
    \end{enumerate}
  \end{enumerate}
\end{defn}

Bu tanıma göre $\mSat{M}{!B\cif !C}[w]$, ancak ve ancak önbileşen~$!B$,
$w$ çevresindeki kürelerin her yerinde yanlışsa ya da $!B$'nin doğru olduğu
ve maddi koşul $!B\lif !C$'nin bu ``$!B$-kabul eden'' küredeki bütün
dünyalarda doğru olduğu bir $S$ küresi varsa geçerlidir. Sezgisel açıklamadan
beklenebileceğinin tersine tanımda $S$'nin $!B$'yi kabul eden \emph{en iç}
küre olmasını istemediğimize dikkat edin. Fakat
\olref{sphere-nonvac}'deki koşul bir $S$ küresi için sağlanıyorsa $S$
içindeki bütün küreler için de sağlanır; dolayısıyla özellikle en içteki küre
için sağlanır.

Ayrıca küre modellerinin tanımı, $!B$'yi kabul eden en içte bir kürenin
\emph{var olmasını} gerektirmez: $!B$'yi kabul eden
$S_1\supsetneq S_2\supsetneq\dots\supsetneq\{w\}$ biçiminde sonsuz bir
küre dizisi bulunabilir ve bu nedenle en içteki $!B$-kabul eden küre
bulunmayabilir. Bu durumda $\mSat{M}{!B\cif !C}[w]$, ancak ve ancak bazı~$i$
için $!B\lif !C$, $S_i,S_{i+1},\dots$ kürelerinin tamamında geçerlidir.

\end{document}
